% cover_letter_dwave.tex -- cover letter for D-Wave Quantum, Senior Scientific
% Software Developer [635], in the design of "Michael Lehn Cover Letter.pages"
% Compile with: lualatex cover_letter_dwave.tex   (or: make)
%
% Layout values are measured from the PDF exported by Pages, in PDF points
% (bp, 1/72 in; TeX's pt is slightly smaller); y coordinates of absolutely
% placed elements are measured from the bottom of the page, as in PDF:
%   A4, body text 62.35bp from the left and right edges, first baseline
%   (salutation) 301.36bp below the top, bottom margin 56.7bp; font Arial;
%   body 10bp with a line pitch of 13.13bp and one empty line between
%   paragraphs; justified with hyphenation and first-fit line breaking.
%   Header block 8bp with an 11bp line pitch at x = 426.38bp.
%
% The phone number is read from private.tex (not committed). The public
% version omits it and is built from this file with
%   lualatex -jobname=cover_letter_dwave_public '\def\PublicVersion{}% cover_letter_dwave.tex -- cover letter for D-Wave Quantum, Senior Scientific
% Software Developer [635], in the design of "Michael Lehn Cover Letter.pages"
% Compile with: lualatex cover_letter_dwave.tex   (or: make)
%
% Layout values are measured from the PDF exported by Pages, in PDF points
% (bp, 1/72 in; TeX's pt is slightly smaller); y coordinates of absolutely
% placed elements are measured from the bottom of the page, as in PDF:
%   A4, body text 62.35bp from the left and right edges, first baseline
%   (salutation) 301.36bp below the top, bottom margin 56.7bp; font Arial;
%   body 10bp with a line pitch of 13.13bp and one empty line between
%   paragraphs; justified with hyphenation and first-fit line breaking.
%   Header block 8bp with an 11bp line pitch at x = 426.38bp.
%
% The phone number is read from private.tex (not committed). The public
% version omits it and is built from this file with
%   lualatex -jobname=cover_letter_dwave_public '\def\PublicVersion{}% cover_letter_dwave.tex -- cover letter for D-Wave Quantum, Senior Scientific
% Software Developer [635], in the design of "Michael Lehn Cover Letter.pages"
% Compile with: lualatex cover_letter_dwave.tex   (or: make)
%
% Layout values are measured from the PDF exported by Pages, in PDF points
% (bp, 1/72 in; TeX's pt is slightly smaller); y coordinates of absolutely
% placed elements are measured from the bottom of the page, as in PDF:
%   A4, body text 62.35bp from the left and right edges, first baseline
%   (salutation) 301.36bp below the top, bottom margin 56.7bp; font Arial;
%   body 10bp with a line pitch of 13.13bp and one empty line between
%   paragraphs; justified with hyphenation and first-fit line breaking.
%   Header block 8bp with an 11bp line pitch at x = 426.38bp.
%
% The phone number is read from private.tex (not committed). The public
% version omits it and is built from this file with
%   lualatex -jobname=cover_letter_dwave_public '\def\PublicVersion{}% cover_letter_dwave.tex -- cover letter for D-Wave Quantum, Senior Scientific
% Software Developer [635], in the design of "Michael Lehn Cover Letter.pages"
% Compile with: lualatex cover_letter_dwave.tex   (or: make)
%
% Layout values are measured from the PDF exported by Pages, in PDF points
% (bp, 1/72 in; TeX's pt is slightly smaller); y coordinates of absolutely
% placed elements are measured from the bottom of the page, as in PDF:
%   A4, body text 62.35bp from the left and right edges, first baseline
%   (salutation) 301.36bp below the top, bottom margin 56.7bp; font Arial;
%   body 10bp with a line pitch of 13.13bp and one empty line between
%   paragraphs; justified with hyphenation and first-fit line breaking.
%   Header block 8bp with an 11bp line pitch at x = 426.38bp.
%
% The phone number is read from private.tex (not committed). The public
% version omits it and is built from this file with
%   lualatex -jobname=cover_letter_dwave_public '\def\PublicVersion{}\input{cover_letter_dwave}'
\documentclass[10pt]{article}

\usepackage[a4paper,left=62.35bp,right=62.35bp,top=291.36bp,bottom=56.7bp,
            headheight=0bp,headsep=0bp,footskip=0bp]{geometry}
\usepackage[american]{babel}
\usepackage{fontspec}
\usepackage[HTML]{xcolor}
\usepackage{graphicx}
\usepackage{eso-pic}
\usepackage{lua-ul}
\usepackage{hyperref}

% Pages breaks lines first-fit; TeX's paragraph optimizer would not
\directlua{dofile(kpse.find_file("firstfit.lua", "lua"))}

% Arial if installed, otherwise its metric-compatible clone Arimo.
% No ligatures, and no TeX input mapping (keeps straight quotes as typed).
\IfFontExistsTF{Arial}
  {\setmainfont{Arial}[Ligatures={CommonOff,TeXOff}]}
  {\setmainfont{Arimo}[Ligatures={CommonOff,TeXOff}]}

\definecolor{link}{HTML}{000099}     % hyperlinks (Pages: RGB 0 0 0.6)

\hypersetup{hidelinks, pdftitle={Michael Lehn -- Cover Letter -- D-Wave Quantum},
            pdfauthor={Michael Lehn}}

% Hyperlink in the Pages style: blue and underlined; the underline is
% 0.586bp thick, centered 0.47bp below the baseline.
\newunderlinetype\linkUnderLine{\leaders\vrule height -0.177bp depth 0.763bp}
\newcommand{\link}[2]{\href{#1}{\textcolor{link}{\linkUnderLine{#2}}}}

% ---------------------------------------------------------------------------
% Letterhead, placed absolutely on the page
\newcommand{\letterdate}{October 29, 2026}

\AddToShipoutPictureBG*{%
  \setlength{\unitlength}{1bp}%
  % Ulm University seal and logo
  \put(56.7,752.9){\includegraphics[width=481.2bp,height=49.4bp]{uulm-logo.png}}%
  % Fold mark and hole-punch mark (105mm and 148mm from the top)
  \put(0,544.3){\color{black}\rule[-0.125bp]{29.6bp}{0.25bp}}%
  \put(0,422.4){\color{black}\rule[-0.125bp]{29.6bp}{0.25bp}}%
  % Sender: first baseline at y = 724.707bp, empty lines between groups
  \put(426.3833,724.707){\vtop{%
    \hsize=165bp \parindent=0bp \parskip=0bp
    \fontsize{8bp}{11bp}\selectfont
    \textbf{Michael Lehn, Ph.D.}\par
    Scientific Software Developer\par
    \strut\par
    Helmholtzstrasse 20\par
    89081 Ulm, Germany\par
    \strut\par
    \ifdefined\PublicVersion\else Mobile:\ \ \PrivatePhone\par\fi
    Email: \link{mailto:michael.lehn@uni-ulm.de}{michael.lehn@uni-ulm.de}\par
    GitHub: \link{https://github.com/michael-lehn}{github.com/michael-lehn}\par
    \strut\par
    \textbf{\letterdate}\par}}%
}

% Private contact details (phone) come from private.tex, which is not
% committed; the private version is not built without it.
\ifdefined\PublicVersion\else
  \InputIfFileExists{private.tex}{}%
    {\errmessage{private.tex not found; it is required for the private version}}
\fi

% ---------------------------------------------------------------------------
\setlength{\parindent}{0bp}
\setlength{\parskip}{13.1346bp}  % one empty line between paragraphs
\setlength{\topskip}{10bp}
\pagestyle{empty}

\begin{document}
\hyphenation{HeraldRT SystemVerilog ulmBLAS FLENS D-Wave}
\fontsize{10bp}{13.1346bp}\selectfont
\frenchspacing       % uniform interword spaces after punctuation, as in Pages
\tolerance=9999      % Pages stretches interword spaces freely in justified text
\emergencystretch=0bp

Dear D-Wave Hiring Team,

I am applying for the Senior Scientific Software Developer position in Burnaby. I am a scientific software developer with over 20 years of hands-on experience in C/C++, numerical mathematics, high-performance computing, and software architecture. At Ulm University, I design, implement, and maintain numerical and systems software together with scientists and engineers. Designing, implementing, debugging, and optimizing scientific software is what I enjoy most and what I want to focus on professionally.

Much of my work has gone into numerical libraries. FLENS, the C++ framework I developed, maps high-level matrix abstractions directly to optimized BLAS and LAPACK kernels without abstraction overhead; my BLAS library ulmBLAS separates portable blocking algorithms from small, architecture-specific micro-kernels and reaches performance close to vendor libraries. Maintaining both over many years has taught me to treat architecture, testing, benchmarking, and documentation as part of the engineering work.

With the Institute of Quantum Technologies at the German Aerospace Center (DLR), I am developing HeraldRT, a SystemVerilog design for FPGA-based real-time control of quantum-memory experiments. It has been simulated, synthesized, placed, and routed; testing on physical hardware is the next step. The work requires translating the physicists' requirements into deterministic hardware behavior, meeting timing constraints, and defining clean hardware/software interfaces; it builds directly on my background in processor architecture, FPGA design, and low-level programming.

As managing director of the Ulm Center for Scientific Computing, I translated problems of industrial and research partners into mathematical models and working software, often implementing the complete application myself. I am used to taking ownership of a problem: understanding the requirements, investigating failures until the cause is clear, implementing a solution, and validating the result.

What attracts me to this position is that it combines hands-on scientific software engineering, from calibration to hardware interfaces and system testing, with direct collaboration with the physicists and hardware engineers who develop D-Wave's quantum processors. The control software I build for quantum-memory experiments at DLR faces similar demands: precise timing, reliable hardware interfaces, and close iteration with the physicists who use it. I am willing to relocate to the Vancouver/Burnaby area.

Thank you for your consideration.

Sincerely,

Michael Lehn
\end{document}
'
\documentclass[10pt]{article}

\usepackage[a4paper,left=62.35bp,right=62.35bp,top=291.36bp,bottom=56.7bp,
            headheight=0bp,headsep=0bp,footskip=0bp]{geometry}
\usepackage[american]{babel}
\usepackage{fontspec}
\usepackage[HTML]{xcolor}
\usepackage{graphicx}
\usepackage{eso-pic}
\usepackage{lua-ul}
\usepackage{hyperref}

% Pages breaks lines first-fit; TeX's paragraph optimizer would not
\directlua{dofile(kpse.find_file("firstfit.lua", "lua"))}

% Arial if installed, otherwise its metric-compatible clone Arimo.
% No ligatures, and no TeX input mapping (keeps straight quotes as typed).
\IfFontExistsTF{Arial}
  {\setmainfont{Arial}[Ligatures={CommonOff,TeXOff}]}
  {\setmainfont{Arimo}[Ligatures={CommonOff,TeXOff}]}

\definecolor{link}{HTML}{000099}     % hyperlinks (Pages: RGB 0 0 0.6)

\hypersetup{hidelinks, pdftitle={Michael Lehn -- Cover Letter -- D-Wave Quantum},
            pdfauthor={Michael Lehn}}

% Hyperlink in the Pages style: blue and underlined; the underline is
% 0.586bp thick, centered 0.47bp below the baseline.
\newunderlinetype\linkUnderLine{\leaders\vrule height -0.177bp depth 0.763bp}
\newcommand{\link}[2]{\href{#1}{\textcolor{link}{\linkUnderLine{#2}}}}

% ---------------------------------------------------------------------------
% Letterhead, placed absolutely on the page
\newcommand{\letterdate}{October 29, 2026}

\AddToShipoutPictureBG*{%
  \setlength{\unitlength}{1bp}%
  % Ulm University seal and logo
  \put(56.7,752.9){\includegraphics[width=481.2bp,height=49.4bp]{uulm-logo.png}}%
  % Fold mark and hole-punch mark (105mm and 148mm from the top)
  \put(0,544.3){\color{black}\rule[-0.125bp]{29.6bp}{0.25bp}}%
  \put(0,422.4){\color{black}\rule[-0.125bp]{29.6bp}{0.25bp}}%
  % Sender: first baseline at y = 724.707bp, empty lines between groups
  \put(426.3833,724.707){\vtop{%
    \hsize=165bp \parindent=0bp \parskip=0bp
    \fontsize{8bp}{11bp}\selectfont
    \textbf{Michael Lehn, Ph.D.}\par
    Scientific Software Developer\par
    \strut\par
    Helmholtzstrasse 20\par
    89081 Ulm, Germany\par
    \strut\par
    \ifdefined\PublicVersion\else Mobile:\ \ \PrivatePhone\par\fi
    Email: \link{mailto:michael.lehn@uni-ulm.de}{michael.lehn@uni-ulm.de}\par
    GitHub: \link{https://github.com/michael-lehn}{github.com/michael-lehn}\par
    \strut\par
    \textbf{\letterdate}\par}}%
}

% Private contact details (phone) come from private.tex, which is not
% committed; the private version is not built without it.
\ifdefined\PublicVersion\else
  \InputIfFileExists{private.tex}{}%
    {\errmessage{private.tex not found; it is required for the private version}}
\fi

% ---------------------------------------------------------------------------
\setlength{\parindent}{0bp}
\setlength{\parskip}{13.1346bp}  % one empty line between paragraphs
\setlength{\topskip}{10bp}
\pagestyle{empty}

\begin{document}
\hyphenation{HeraldRT SystemVerilog ulmBLAS FLENS D-Wave}
\fontsize{10bp}{13.1346bp}\selectfont
\frenchspacing       % uniform interword spaces after punctuation, as in Pages
\tolerance=9999      % Pages stretches interword spaces freely in justified text
\emergencystretch=0bp

Dear D-Wave Hiring Team,

I am applying for the Senior Scientific Software Developer position in Burnaby. I am a scientific software developer with over 20 years of hands-on experience in C/C++, numerical mathematics, high-performance computing, and software architecture. At Ulm University, I design, implement, and maintain numerical and systems software together with scientists and engineers. Designing, implementing, debugging, and optimizing scientific software is what I enjoy most and what I want to focus on professionally.

Much of my work has gone into numerical libraries. FLENS, the C++ framework I developed, maps high-level matrix abstractions directly to optimized BLAS and LAPACK kernels without abstraction overhead; my BLAS library ulmBLAS separates portable blocking algorithms from small, architecture-specific micro-kernels and reaches performance close to vendor libraries. Maintaining both over many years has taught me to treat architecture, testing, benchmarking, and documentation as part of the engineering work.

With the Institute of Quantum Technologies at the German Aerospace Center (DLR), I am developing HeraldRT, a SystemVerilog design for FPGA-based real-time control of quantum-memory experiments. It has been simulated, synthesized, placed, and routed; testing on physical hardware is the next step. The work requires translating the physicists' requirements into deterministic hardware behavior, meeting timing constraints, and defining clean hardware/software interfaces; it builds directly on my background in processor architecture, FPGA design, and low-level programming.

As managing director of the Ulm Center for Scientific Computing, I translated problems of industrial and research partners into mathematical models and working software, often implementing the complete application myself. I am used to taking ownership of a problem: understanding the requirements, investigating failures until the cause is clear, implementing a solution, and validating the result.

What attracts me to this position is that it combines hands-on scientific software engineering, from calibration to hardware interfaces and system testing, with direct collaboration with the physicists and hardware engineers who develop D-Wave's quantum processors. The control software I build for quantum-memory experiments at DLR faces similar demands: precise timing, reliable hardware interfaces, and close iteration with the physicists who use it. I am willing to relocate to the Vancouver/Burnaby area.

Thank you for your consideration.

Sincerely,

Michael Lehn
\end{document}
'
\documentclass[10pt]{article}

\usepackage[a4paper,left=62.35bp,right=62.35bp,top=291.36bp,bottom=56.7bp,
            headheight=0bp,headsep=0bp,footskip=0bp]{geometry}
\usepackage[american]{babel}
\usepackage{fontspec}
\usepackage[HTML]{xcolor}
\usepackage{graphicx}
\usepackage{eso-pic}
\usepackage{lua-ul}
\usepackage{hyperref}

% Pages breaks lines first-fit; TeX's paragraph optimizer would not
\directlua{dofile(kpse.find_file("firstfit.lua", "lua"))}

% Arial if installed, otherwise its metric-compatible clone Arimo.
% No ligatures, and no TeX input mapping (keeps straight quotes as typed).
\IfFontExistsTF{Arial}
  {\setmainfont{Arial}[Ligatures={CommonOff,TeXOff}]}
  {\setmainfont{Arimo}[Ligatures={CommonOff,TeXOff}]}

\definecolor{link}{HTML}{000099}     % hyperlinks (Pages: RGB 0 0 0.6)

\hypersetup{hidelinks, pdftitle={Michael Lehn -- Cover Letter -- D-Wave Quantum},
            pdfauthor={Michael Lehn}}

% Hyperlink in the Pages style: blue and underlined; the underline is
% 0.586bp thick, centered 0.47bp below the baseline.
\newunderlinetype\linkUnderLine{\leaders\vrule height -0.177bp depth 0.763bp}
\newcommand{\link}[2]{\href{#1}{\textcolor{link}{\linkUnderLine{#2}}}}

% ---------------------------------------------------------------------------
% Letterhead, placed absolutely on the page
\newcommand{\letterdate}{October 29, 2026}

\AddToShipoutPictureBG*{%
  \setlength{\unitlength}{1bp}%
  % Ulm University seal and logo
  \put(56.7,752.9){\includegraphics[width=481.2bp,height=49.4bp]{uulm-logo.png}}%
  % Fold mark and hole-punch mark (105mm and 148mm from the top)
  \put(0,544.3){\color{black}\rule[-0.125bp]{29.6bp}{0.25bp}}%
  \put(0,422.4){\color{black}\rule[-0.125bp]{29.6bp}{0.25bp}}%
  % Sender: first baseline at y = 724.707bp, empty lines between groups
  \put(426.3833,724.707){\vtop{%
    \hsize=165bp \parindent=0bp \parskip=0bp
    \fontsize{8bp}{11bp}\selectfont
    \textbf{Michael Lehn, Ph.D.}\par
    Scientific Software Developer\par
    \strut\par
    Helmholtzstrasse 20\par
    89081 Ulm, Germany\par
    \strut\par
    \ifdefined\PublicVersion\else Mobile:\ \ \PrivatePhone\par\fi
    Email: \link{mailto:michael.lehn@uni-ulm.de}{michael.lehn@uni-ulm.de}\par
    GitHub: \link{https://github.com/michael-lehn}{github.com/michael-lehn}\par
    \strut\par
    \textbf{\letterdate}\par}}%
}

% Private contact details (phone) come from private.tex, which is not
% committed; the private version is not built without it.
\ifdefined\PublicVersion\else
  \InputIfFileExists{private.tex}{}%
    {\errmessage{private.tex not found; it is required for the private version}}
\fi

% ---------------------------------------------------------------------------
\setlength{\parindent}{0bp}
\setlength{\parskip}{13.1346bp}  % one empty line between paragraphs
\setlength{\topskip}{10bp}
\pagestyle{empty}

\begin{document}
\hyphenation{HeraldRT SystemVerilog ulmBLAS FLENS D-Wave}
\fontsize{10bp}{13.1346bp}\selectfont
\frenchspacing       % uniform interword spaces after punctuation, as in Pages
\tolerance=9999      % Pages stretches interword spaces freely in justified text
\emergencystretch=0bp

Dear D-Wave Hiring Team,

I am applying for the Senior Scientific Software Developer position in Burnaby. I am a scientific software developer with over 20 years of hands-on experience in C/C++, numerical mathematics, high-performance computing, and software architecture. At Ulm University, I design, implement, and maintain numerical and systems software together with scientists and engineers. Designing, implementing, debugging, and optimizing scientific software is what I enjoy most and what I want to focus on professionally.

Much of my work has gone into numerical libraries. FLENS, the C++ framework I developed, maps high-level matrix abstractions directly to optimized BLAS and LAPACK kernels without abstraction overhead; my BLAS library ulmBLAS separates portable blocking algorithms from small, architecture-specific micro-kernels and reaches performance close to vendor libraries. Maintaining both over many years has taught me to treat architecture, testing, benchmarking, and documentation as part of the engineering work.

With the Institute of Quantum Technologies at the German Aerospace Center (DLR), I am developing HeraldRT, a SystemVerilog design for FPGA-based real-time control of quantum-memory experiments. It has been simulated, synthesized, placed, and routed; testing on physical hardware is the next step. The work requires translating the physicists' requirements into deterministic hardware behavior, meeting timing constraints, and defining clean hardware/software interfaces; it builds directly on my background in processor architecture, FPGA design, and low-level programming.

As managing director of the Ulm Center for Scientific Computing, I translated problems of industrial and research partners into mathematical models and working software, often implementing the complete application myself. I am used to taking ownership of a problem: understanding the requirements, investigating failures until the cause is clear, implementing a solution, and validating the result.

What attracts me to this position is that it combines hands-on scientific software engineering, from calibration to hardware interfaces and system testing, with direct collaboration with the physicists and hardware engineers who develop D-Wave's quantum processors. The control software I build for quantum-memory experiments at DLR faces similar demands: precise timing, reliable hardware interfaces, and close iteration with the physicists who use it. I am willing to relocate to the Vancouver/Burnaby area.

Thank you for your consideration.

Sincerely,

Michael Lehn
\end{document}
'
\documentclass[10pt]{article}

\usepackage[a4paper,left=62.35bp,right=62.35bp,top=291.36bp,bottom=56.7bp,
            headheight=0bp,headsep=0bp,footskip=0bp]{geometry}
\usepackage[american]{babel}
\usepackage{fontspec}
\usepackage[HTML]{xcolor}
\usepackage{graphicx}
\usepackage{eso-pic}
\usepackage{lua-ul}
\usepackage{hyperref}

% Pages breaks lines first-fit; TeX's paragraph optimizer would not
\directlua{dofile(kpse.find_file("firstfit.lua", "lua"))}

% Arial if installed, otherwise its metric-compatible clone Arimo.
% No ligatures, and no TeX input mapping (keeps straight quotes as typed).
\IfFontExistsTF{Arial}
  {\setmainfont{Arial}[Ligatures={CommonOff,TeXOff}]}
  {\setmainfont{Arimo}[Ligatures={CommonOff,TeXOff}]}

\definecolor{link}{HTML}{000099}     % hyperlinks (Pages: RGB 0 0 0.6)

\hypersetup{hidelinks, pdftitle={Michael Lehn -- Cover Letter -- D-Wave Quantum},
            pdfauthor={Michael Lehn}}

% Hyperlink in the Pages style: blue and underlined; the underline is
% 0.586bp thick, centered 0.47bp below the baseline.
\newunderlinetype\linkUnderLine{\leaders\vrule height -0.177bp depth 0.763bp}
\newcommand{\link}[2]{\href{#1}{\textcolor{link}{\linkUnderLine{#2}}}}

% ---------------------------------------------------------------------------
% Letterhead, placed absolutely on the page
\newcommand{\letterdate}{October 29, 2026}

\AddToShipoutPictureBG*{%
  \setlength{\unitlength}{1bp}%
  % Ulm University seal and logo
  \put(56.7,752.9){\includegraphics[width=481.2bp,height=49.4bp]{uulm-logo.png}}%
  % Fold mark and hole-punch mark (105mm and 148mm from the top)
  \put(0,544.3){\color{black}\rule[-0.125bp]{29.6bp}{0.25bp}}%
  \put(0,422.4){\color{black}\rule[-0.125bp]{29.6bp}{0.25bp}}%
  % Sender: first baseline at y = 724.707bp, empty lines between groups
  \put(426.3833,724.707){\vtop{%
    \hsize=165bp \parindent=0bp \parskip=0bp
    \fontsize{8bp}{11bp}\selectfont
    \textbf{Michael Lehn, Ph.D.}\par
    Scientific Software Developer\par
    \strut\par
    Helmholtzstrasse 20\par
    89081 Ulm, Germany\par
    \strut\par
    \ifdefined\PublicVersion\else Mobile:\ \ \PrivatePhone\par\fi
    Email: \link{mailto:michael.lehn@uni-ulm.de}{michael.lehn@uni-ulm.de}\par
    GitHub: \link{https://github.com/michael-lehn}{github.com/michael-lehn}\par
    \strut\par
    \textbf{\letterdate}\par}}%
}

% Private contact details (phone) come from private.tex, which is not
% committed; the private version is not built without it.
\ifdefined\PublicVersion\else
  \InputIfFileExists{private.tex}{}%
    {\errmessage{private.tex not found; it is required for the private version}}
\fi

% ---------------------------------------------------------------------------
\setlength{\parindent}{0bp}
\setlength{\parskip}{13.1346bp}  % one empty line between paragraphs
\setlength{\topskip}{10bp}
\pagestyle{empty}

\begin{document}
\hyphenation{HeraldRT SystemVerilog ulmBLAS FLENS D-Wave}
\fontsize{10bp}{13.1346bp}\selectfont
\frenchspacing       % uniform interword spaces after punctuation, as in Pages
\tolerance=9999      % Pages stretches interword spaces freely in justified text
\emergencystretch=0bp

Dear D-Wave Hiring Team,

I am applying for the Senior Scientific Software Developer position in Burnaby. I am a scientific software developer with over 20 years of hands-on experience in C/C++, numerical mathematics, high-performance computing, and software architecture. At Ulm University, I design, implement, and maintain numerical and systems software together with scientists and engineers. Designing, implementing, debugging, and optimizing scientific software is what I enjoy most and what I want to focus on professionally.

Much of my work has gone into numerical libraries. FLENS, the C++ framework I developed, maps high-level matrix abstractions directly to optimized BLAS and LAPACK kernels without abstraction overhead; my BLAS library ulmBLAS separates portable blocking algorithms from small, architecture-specific micro-kernels and reaches performance close to vendor libraries. Maintaining both over many years has taught me to treat architecture, testing, benchmarking, and documentation as part of the engineering work.

With the Institute of Quantum Technologies at the German Aerospace Center (DLR), I am developing HeraldRT, a SystemVerilog design for FPGA-based real-time control of quantum-memory experiments. It has been simulated, synthesized, placed, and routed; testing on physical hardware is the next step. The work requires translating the physicists' requirements into deterministic hardware behavior, meeting timing constraints, and defining clean hardware/software interfaces; it builds directly on my background in processor architecture, FPGA design, and low-level programming.

As managing director of the Ulm Center for Scientific Computing, I translated problems of industrial and research partners into mathematical models and working software, often implementing the complete application myself. I am used to taking ownership of a problem: understanding the requirements, investigating failures until the cause is clear, implementing a solution, and validating the result.

What attracts me to this position is that it combines hands-on scientific software engineering, from calibration to hardware interfaces and system testing, with direct collaboration with the physicists and hardware engineers who develop D-Wave's quantum processors. The control software I build for quantum-memory experiments at DLR faces similar demands: precise timing, reliable hardware interfaces, and close iteration with the physicists who use it. I am willing to relocate to the Vancouver/Burnaby area.

Thank you for your consideration.

Sincerely,

Michael Lehn
\end{document}
